% visualStyle.tex — Single source of truth for every color used in the
% thesis: tables, the Gantt chart, and TikZ diagrams (images/*.tex).
%
% Palette: Okabe & Ito (2008), "Color Universal Design" — the standard
% categorical palette distinguishable under protanopia, deuteranopia and
% tritanopia. Every color in the thesis must come from here; do not
% \definecolor or use raw color names (fill=orange, fill=green, fill=red, ...)
% ad hoc inside chapters/, tables/*.tex or images/*.tex. Raw xcolor names in
% particular (plain "red"/"green"/"yellow") are NOT CVD-safe — red/green is
% the single most common color-blindness confusion pair.

% ── Okabe-Ito base palette ────────────────────────────────────────────────
\definecolor{cbOrange}{RGB}{230,159,0}
\definecolor{cbSkyBlue}{RGB}{86,180,233}
\definecolor{cbGreen}{RGB}{0,158,115}     % bluish green
\definecolor{cbYellow}{RGB}{240,228,66}
\definecolor{cbBlue}{RGB}{0,114,178}
\definecolor{cbVermillion}{RGB}{213,94,0}
\definecolor{cbPurple}{RGB}{204,121,167}
\definecolor{cbGrey}{RGB}{166,166,166}

% ── Semantic roles (used by every table in the thesis) ───────────────────
% Tint levels below are chosen against two constraints, not eyeballed:
%  1. WCAG 1.4.11 (non-text contrast, >=3:1) / 1.4.3 (text contrast, >=4.5:1)
%     — verified against black text for the band colors below (~9-13:1,
%     comfortably AAA) and against white text for the solid header (~5.2:1,
%     AA). A pale tint (e.g. !35 or lower) still clears contrast but reads
%     as washed-out/grey on paper — the "too dull" failure mode — because
%     contrast alone doesn't capture perceived saturation/vividness.
%  2. Okabe & Ito (2008) hue choices are kept exactly (blue/orange is the
%     most CVD-robust 2-hue pair); only the tint (mix-with-white) percentage
%     is raised for a visibly saturated, "colorful but not extravagant" look.
%
% Header row: solid (untinted) blue background, white bold text.
\colorlet{tableHeaderColor}{cbBlue}
% Zebra striping on long/data-heavy tables: a visible-but-subtle tint of the
% header hue — enough to guide the eye across a row without competing with
% the semantic band colors below. Kept lighter than those on purpose: zebra
% striping is a navigation aid, not information-bearing.
\colorlet{tableAltRowColor}{cbBlue!12}
% Wavelet coefficient examples (tables/{regular,packet}WaveletExample.tex):
% grey = raw signal label, sky blue = row/column index label,
% blue = approximation band (cA), orange = detail band (cD).
% Blue/orange is the single most robust 2-hue pair across all CVD types.
% These ARE information-bearing (which band a cell belongs to), so they are
% tinted much more strongly than the zebra stripe above.
\colorlet{tableSignalLabelColor}{cbGrey}
\colorlet{tableIndexLabelColor}{cbSkyBlue}
\colorlet{tableSignalValuesColor}{cbGreen!55}
\colorlet{tableApproxBandColor}{cbBlue!55}
\colorlet{tableDetailBandColor}{cbOrange!55}

% Gantt chart (chapters/08-proposedApproach.tex, tab:cronogramaGantt):
% same palette, kept here so the chart matches the tables it sits beside.
\colorlet{ganttBarColor}{cbGreen}
\colorlet{ganttMilestoneColor}{cbGreen}
\colorlet{ganttTitleColor}{cbBlue!45}
\colorlet{ganttHighlightColor}{cbOrange}

% ── Reusable header row ────────────────────────────────────────────────────
% Usage inside a tabular/longtable, as the first thing in a header row:
%   \tableHeaderRow \textbf{A} & \textbf{B} \\
\newcommand{\tableHeaderRow}{\rowcolor{tableHeaderColor}\color{white}}

% ── Zebra striping for table/longtable bodies ─────────────────────────────
% Call once, right before \begin{tabular}/\begin{longtable}. Colors rows
% 2, 4, 6, ... (i.e. every data row after a 1-row header) with
% tableAltRowColor; row 1 (the header) is left untouched so \tableHeaderRow
% still controls it. NOTE: \rowcolor itself must always be the literal first
% token of a row — never wrap it in another macro (e.g. a counter-stepper),
% since colortbl's \noalign trick requires nothing to precede it.
\newcommand{\tableZebraStripe}{\rowcolors{2}{white}{tableAltRowColor}}

% ── Wide/wrapping text column, vertically centered ────────────────────────
% Plain p{<width>} columns top-align their content, while ordinary c/l/r
% columns baseline/vertically-center theirs — invisible in a plain \hline
% table, but a visible split (header text sitting below the colored fill)
% once a \rowcolor-filled header row mixes p{} with c/l/r columns. Use the
% column type Y{<width>} instead of p{<width>} in any column spec that will
% share a row with \tableHeaderRow, e.g. \begin{tabular}{@{}Y{10cm}cc@{}}.
\newcolumntype{Y}[1]{>{\raggedright\arraybackslash}m{#1}}

% ── Diagram / flowchart semantic roles (images/*.tex) ─────────────────────
% The thesis has ~20 TikZ flowcharts and node diagrams that all reuse the
% same handful of roles (start/end bubble, process box, decision diamond,
% arrows, autoencoder layer nodes). They used to \definecolor / fill= with
% raw xcolor names (orange, green, yellow, red, gray) independently per
% file — not CVD-safe (plain red/green/yellow is the classic confusion
% triple) and impossible to re-theme consistently. Every images/*.tex file
% must reference these instead of a raw color name.
\colorlet{flowStartEndColor}{cbOrange}
\colorlet{flowProcessColor}{cbGreen}
\colorlet{flowDecisionColor}{cbYellow}
\colorlet{flowArrowColor}{cbGrey}
% Autoencoder-style node diagrams (images/{autoencoder,denoisingAutoencoder,residuo}*.tex):
% grey = input, yellow = hidden/code/latent layer, vermillion = output/reconstruction.
% (Vermillion replaces the original "red" — true red is excluded from the
% Okabe-Ito set for the same red/green confusion reason.)
\colorlet{diagramInputColor}{cbGrey}
\colorlet{diagramHiddenColor}{cbYellow}
\colorlet{diagramOutputColor}{cbVermillion}
% Intermediate/corrupted stage between input and hidden layer (e.g. the
% noised x_hat in images/denoisingAutoencoder.tex, or the pre-activation
% sum in images/residuo.tex) — kept visually distinct from all three roles
% above.
\colorlet{diagramIntermediateColor}{cbSkyBlue}
% Tree diagrams (images/afasias.tex): darker tints so white text keeps
% AA contrast on non-header-sized diagram text.
\colorlet{diagramBranchColor}{cbGreen!80!black}
\colorlet{diagramLeafColor}{cbBlue!80!black}
% Feature-vector pipeline diagrams (images/{bark,mel,linear,lfcc,bfcc}...Vect.tex,
% images/{mel,linear}BandEnergies.tex): single accent color for box/ellipse
% outlines. Okabe-Ito blue, darkened for outline contrast against white fill.
\colorlet{diagramOutlineColor}{cbBlue!80!black}

% ── Marcação temporária de revisão (REMOVER antes da versão final) ────────
% Destaca um trecho que ainda depende de trabalho pendente, para que não
% passe despercebido em uma releitura. Usa cbYellow da paleta Okabe-Ito (e
% não o "yellow" bruto do xcolor, que não é CVD-safe); o tom é rebaixado
% para !40 porque o texto impresso por cima é preto e precisa manter
% contraste AA.
%
% \pendenteMarca{...} destaca um trecho CURTO e sem matemática: usa
% \colorbox, que não quebra linha nem sobrevive a $...$ — por isso não deve
% envolver parágrafos inteiros nem fórmulas. Para blocos, use \pendenteNota,
% que colore apenas o rótulo e deixa o corpo do texto em itálico comum.
\colorlet{pendenteColor}{cbYellow!40}
\newcommand{\pendenteMarca}[1]{\colorbox{pendenteColor}{#1}}
\newcommand{\pendenteNota}[1]{%
	\par\noindent\colorbox{pendenteColor}{\textbf{[PENDENTE]}}\hspace{0.4em}\textit{#1}\par}
